\documentclass[11pt,a4paper]{article}
\usepackage[T1]{fontenc}
\usepackage[utf8]{inputenc}
\usepackage[english]{babel}
\usepackage{lmodern}
\usepackage[margin=2.3cm]{geometry}
\usepackage{amsmath,amssymb,mathtools,stmaryrd}
\usepackage{booktabs,longtable,array}
\usepackage[most]{tcolorbox}
\usepackage[hidelinks]{hyperref}

% ------------------------------------------------------------------
% Notation
\newcommand{\limp}{\Rightarrow}
\newcommand{\nat}{\mathbb{N}}
\newcommand{\intg}{\mathbb{Z}}
\newcommand{\pfun}{\mathrel{+\!\!\!\!\!\rightarrow}}
\newcommand{\dom}{\mathrm{dom}}
\newcommand{\WD}{\mathit{WD}}
\newcommand{\BA}{\mathit{BA}}
\newcommand{\AX}{\mathit{AX}(s,c)}
\newcommand{\Th}{\mathit{Th}}
\newcommand{\sff}[1]{\textsf{#1}}
\newcommand{\po}[1]{\textsf{\small #1}}
\newcommand{\trm}{\mathsf{trm}}
\newcommand{\Fis}{\mathsf{Fis}}
\newcommand{\fis}{\mathsf{fis}}
\newcommand{\ligne}[1]{(l.~#1)}

% Box describing a generation rule
\newtcolorbox{regle}[2][]{enhanced,breakable,colback=blue!3,colframe=blue!45!black,
  fonttitle=\bfseries,title={#2},#1}
\newtcolorbox{remarque}[1][]{enhanced,breakable,colback=orange!5,colframe=orange!60!black,
  fonttitle=\bfseries,title={Remark},#1}
\newtcolorbox{rodin}[1][]{enhanced,breakable,colback=green!4,colframe=green!40!black,
  fonttitle=\bfseries,title={In Rodin: limitations and adjustments},#1}

\setlength{\emergencystretch}{3em}
\newcommand{\champ}[1]{\par\smallskip\noindent\textbf{#1}\enspace}

\title{Generation rules for proof obligations\\
  in the chapter ``The Event-B Modelling Language''\\[2pt]
  \large Analysis and precise description}
\author{Dominique Méry}
\date{\today}
\begin{document}
\bibliographystyle{plain}
\maketitle

\begin{abstract}
  This document gives a precise description of the rules for generating the
  verification conditions (\emph{proof obligations}, written POs below) presented in
  the Event-B chapter \cite{buneventb}  of the book \emph{Modelling Software-based Systems:
  Correct by Construction Paradigm}\cite{bun}. For each rule, it gives its origin
  in the text, the model element that triggers it, the exact form of the sequent it
  produces, its naming convention (Rodin style), its justification and, when the
  chapter discusses them, the limitations and adjustments of Rodin. 
 \end{abstract}

\tableofcontents

% ==================================================================
\section{Purpose, scope and method}

The chapter under analysis (compiled version \texttt{main-C1.pdf}) presents the POs
in three stages:
\begin{enumerate}
  \item \textbf{semantic foundations} (§1.2): transition system, induction
    principle, fixed-point characterisation, equivalence of the \emph{wp} and
    relational styles;
  \item the \textbf{POs of a context and of an abstract machine} (§1.3);
  \item the \textbf{POs of refinement} (§1.4), complemented by event merging and by
    a semantic discussion of deadlock freedom.
\end{enumerate}
Section~1.5 (abacus, ``Rule 1'' and ``Rule 2'') describes the rules for using an
abacus and defines no new PO generation rule. It illustrates how the rules
described here are applied, and is therefore outside the scope of this document.

The text always follows the same generation method:
\begin{description}
  \item[(M1) Global condition.] Start from a global semantic condition
    (preservation of an invariant, refinement in \emph{wp}, etc.).
  \item[(M2) Sequent form.] Rewrite it as a sequent by removing the universal
    quantifications over the state, which become free variables.
  \item[(M3) Decomposition.] Apply two sequent-calculus rules
    (§\ref{sec:simplif}) to obtain elementary sequents.
  \item[(M4) Elimination of existentials.] When possible, replace an existential
    quantification in the conclusion by a \emph{witness} supplied by the modeller,
    which gives rise to a PO for the feasibility of the witness.
  \item[(M5) Naming.] Each elementary sequent is named after the label
    $\ell(\mathit{text})$ of the formal text (axiom, invariant, guard, action,
    theorem) that produced it, in the form \po{element/label/TYPE}.
\end{description}

% ==================================================================
\section{Framework, notation and basic rules}

\subsection{Event-based model}
An event-based model of a system \textsc{S} (def.~\emph{event-based model}) is a
quintuple
\[
  \mathit{EM}(S) = (\AX,\ \{x\},\ \mathcal{V}\!\mathit{alues},\ \mathit{Init}(x),\ \{e_0,\dots,e_n\})
\]
where $s$ are the basic sets, $c$ the constants, $\AX$ the axiomatic theory (the
context), $x$ the flexible variables, $\mathit{Init}$ the initial predicate and
$e_i$ the events. The event $e_0$ is \sff{skip}. The chapter now writes states
$\sigma,\sigma'$, the letter $s$ being reserved for the basic sets,
and writes the refinement events in sans-serif type, $\mathsf{a}$ and $\mathsf{c}$,
to distinguish the concrete event $\mathsf{c}$ from the list $c$ of constants.

\subsection{Normalised form of an event}\label{sec:forme}
An event is written
\[
  \sff{e} \;\widehat{=}\; \textbf{any } u \textbf{ where } G(u,s,c,x)
  \textbf{ then } x :\mid \mathit{BAP}(u,s,c,x,x') \textbf{ end}
\]
and its before-after relation is
\[
  \BA(\sff{e})(s,c,x,x') \;\widehat{=}\; \exists u\cdot G(u,s,c,x) \wedge \mathit{BAP}(u,s,c,x,x').
\]
The global transition relation is $\mathit{Next} \widehat{=} \bigvee_{i} \BA(e_i)$,
and $\mathit{Next}^\star$ is its reflexive transitive closure, defined as the
\emph{least} relation satisfying the recursive equation of the chapter.

\subsection{Notation used in the sequents}
\begin{center}
\begin{tabular}{>{$}l<{$} p{0.72\linewidth}}
\toprule
\AX & conjunction of the axioms of the context that is seen (\sff{sees}) and of the contexts it extends\\
\Th(s,c) & theorems of the context that precede the theorem under consideration\\
\Th_i(s,c) & theorems located \emph{before} the machine theorem $th_i$ (for $i=1$, only those of the context)\\
I(s,c,x) & invariant $I_1\wedge\dots\wedge I_r$ of the machine, labelled $inv_1,\dots,inv_r$\\
\mathit{SAFE}_i(s,c,x) & safety properties (clause \sff{THEOREMS}), $i\in 1..h$; $A = \bigwedge_i \mathit{SAFE}_i$\\
\mathit{initBAP} & before-after predicate of the initialisation, split by actions $act_i$; $\mathit{Init}(s,c,x)\equiv\mathit{initBAP}(s,c,x)$\\
G,\ \mathit{BAP} & guard and before-after predicate of the abstract event\\
J(s,c,x,y) & gluing invariant $J_1\wedge\dots\wedge J_r$, labelled $jnv_i$, linking $x$ (abstract) and $y$ (concrete)\\
H,\ \mathit{CBAP} & guard and before-after predicate of the concrete event (parameters $v$)\\
\mathit{WP}(u,s,c,v,y) & witness for the abstract parameter $u$ (\emph{Witness Parameter})\\
\mathit{WV}(v,s,c,y',x') & witness for the abstract after-value $x'$ (\emph{Witness Variable}); at initialisation, $\mathit{WV}(s,c,y',x')$\\
\mathit{varexp} & the single variant of the machine (natural number or finite set)\\
\bottomrule
\end{tabular}
\end{center}

\subsection{Simplification of typing}
The semantic conditions are first written in the form
$\AX \vdash \forall x\in\mathcal{V}\!\mathit{alues}\cdot P(x)$. Since $x$ has no free
occurrence in $\AX$, they are rewritten as $\AX, x\in\mathcal{V}\!\mathit{alues}\vdash P(x)$.
The typing hypothesis $x\in\mathcal{V}\!\mathit{alues}$ is then assumed to be carried
by the invariant and the initialisation, and is omitted.

\subsection{The two sequent decomposition rules}\label{sec:simplif}
The whole generation rests on two rules (property ``simplification of sequents''
and property ``Introduction/Elimination of connector $\wedge$''):
\begin{regle}{Rule S1: implication introduction}
  \[
    \AX, \mathcal{H} \vdash P_1 \limp P_2 \limp \dots \limp P_n \limp Q
    \quad\text{is equivalent to}\quad
    \AX, \mathcal{H}, P_1, \dots, P_n \vdash Q .
  \]
\end{regle}
\begin{regle}{Rule S2: splitting a conjunction}
  \[
    \Gamma \vdash \textstyle\bigwedge_{i\in 1..t} P_i
    \quad\text{is equivalent to}\quad
    \Gamma \vdash P_i \ \text{ for all } i\in 1..t .
  \]
  This is the rule that produces \emph{one PO per label}: one per invariant, per
  theorem, per action, etc.
\end{regle}
Two implicit rules complete these explicit ones. The text uses them in the
derivations without naming them:
\begin{itemize}
  \item \textbf{S3, existential in a hypothesis}:
    $\Gamma, \exists v\cdot P(v) \vdash Q$ is equivalent to $\Gamma, P(v)\vdash Q$
    if $v$ is not free in $\Gamma, Q$ (the variable becomes a free parameter of the
    sequent);
  \item \textbf{S4, existential in the conclusion by a witness}: if
    $\Gamma, W(u) \vdash P(u)$ and $\Gamma\vdash\exists u\cdot W(u)$, then
    $\Gamma\vdash\exists u\cdot P(u)$. This is the basis of the \po{WFIS} POs
    (§\ref{sec:wfis}); the second premise is essential.
\end{itemize}

\subsection{Naming convention}
POs are named as in Rodin: \po{object/label/TYPE}, where \emph{object} is an event
(or \po{INITIALISATION}), a theorem or an axiom, \emph{label} is the label of the
invariant, action, guard or witness concerned, and \emph{TYPE} is one of the
following suffixes: \po{WD}, \po{THM}, \po{INV}, \po{FIS}, \po{GRD}, \po{SIM},
\po{WFIS}, \po{WWD}, \po{NAT}, \po{FIN}, \po{VAR}, \po{MRG}, \po{DLF}. For variants,
Rodin inserts no label: \po{e/NAT}, \po{e/FIN}, \po{e/VAR}. The suffixes \po{-WIT},
\po{-P} and \po{-V} (\po{GRD-WIT}, \po{WFIS-P}, \dots) belong to the chapter: Rodin
produces \po{GRD} and \po{SIM} with the witness predicates as hypotheses.

% ==================================================================
\section{Foundations: where the POs come from}

\subsection{Induction principle}
\begin{regle}{Induction principle (property \emph{induction principle})}
  $P(x)$ is a safety property of \textsc{S} if, and only if, there exists $I(x)$
  such that, for all $x,x'$:
  \[
    \text{(1) } \mathit{Init}(x)\limp I(x) \qquad
    \text{(2) } I(x)\limp P(x) \qquad
    \text{(3) } I(x)\wedge \mathit{Next}(s,c,x,x') \limp I(x').
  \]
\end{regle}
The \emph{equivalence of induction principles} replaces (3) by one condition per
event: $\forall i\cdot I(x)\wedge\BA(e_i)(s,c,x,x')\limp I(x')$. This is rule S2
applied to the disjunction in the hypothesis, since
$(\bigvee_i A_i)\limp B \equiv \bigwedge_i (A_i\limp B)$. This yields the three
families of POs of the definition ``Verification conditions for a system
\textsc{S}'' (where the safety property is written $A(x)$):
\[
  \AX\vdash \mathit{Init}\limp I, \qquad
  \AX\vdash I\wedge \BA(\sff{e})\limp I' \ (\forall \sff{e}), \qquad
  \AX\vdash I\limp A .
\]

\subsection{Fixed-point characterisation}\label{sec:pf}
Let $F = \lambda X\cdot \mathit{INIT}\cup\mathit{Next}[X]$ on the lattice
$(\mathcal{P}(\mathcal{V}\!\mathit{alues}),\subseteq)$. The \emph{inductive
invariants} are exactly the sets $I$ such that $F(I)\subseteq I$ (pre-fixed points),
that is, $\mathit{INIT}\subseteq I$ and $\mathit{Next}[I]\subseteq I$; they are not
only the solutions of $X=F(X)$. The least fixed point
$\mathit{lfp}(F)$, the set of reachable states, is their intersection
(Knaster--Tarski): every inductive invariant \emph{contains} the reachable states.
The POs therefore only check that a \emph{proposed} invariant is indeed inductive:
this is the correct-by-construction approach.

\subsection{Termination and feasibility of an event}
For an event, $\trm(\sff{e}) \widehat{=} [\sff{e}]\mathit{true}$ is always
$\mathit{true}$ (under the WD conditions): an event always terminates. Only its
feasibility is at stake. The chapter distinguishes two predicates:
\begin{align*}
  \fis(\sff{e})(s,c,x) &\equiv \exists u\cdot\big(G(u,s,c,x)\wedge\exists x'\cdot \mathit{BAP}(u,s,c,x,x')\big)
    && \text{(global feasibility: the actual guard)},\\
  \Fis(\sff{e})(s,c,x) &\widehat{=} \forall u\cdot\big(G(u,s,c,x)\limp \exists x'\cdot \mathit{BAP}(u,s,c,x,x')\big)
    && \text{(per-parameter feasibility)}.
\end{align*}
The obligation \po{FIS} requires $I\limp\Fis(\sff{e})$: it rules out the
\emph{hidden miracle}, a guard that holds while the action produces no state. The
name \po{FIS} is reserved for the proof obligation.

\subsection{Equivalence of the \emph{wp} and relational styles}
For an event \sff{e} and a property $P$:
\[
  [\sff{e}]P(s,c,x) \equiv \forall u\cdot\big(G(u,s,c,x)\limp \forall x'\cdot(\mathit{BAP}(u,s,c,x,x')\limp P(s,c,x'))\big).
\]
The \emph{wp} condition (Atelier~B style)
$\vdash I\limp\Fis(\sff{e})\wedge[\sff{e}]I$ is equivalent to the pair of
relational conditions (Rodin style), free variables being universally quantified
(property ``\emph{wp} and relational styles'' of the chapter, with its proof):
\[
  \text{(FIS)}\ \ I\wedge G \limp \exists x'\cdot \mathit{BAP},
  \qquad
  \text{(INV)}\ \ I\wedge G\wedge \mathit{BAP}\limp I(x').
\]
\textbf{Consequence:} invariant preservation always splits into an \emph{induction
step} (INV) and a \emph{feasibility proof} (FIS). The transformer $[\sff{e}]$ alone
says nothing about feasibility: a parameter value with no next state satisfies
every postcondition. This is why $\Fis(\sff{e})$ appears in the
\emph{conclusion} of the \emph{wp} condition, not as a hypothesis. The chapter's
definition ``verification condition \sff{e}'' (def.~\texttt{def:poe}) is
$\AX\vdash I\limp\Fis(\sff{e})\wedge[\sff{e}]I$.

% ==================================================================
\section{Rules for a context}

\begin{regle}{R1 \quad \po{axm/WD}, \po{th/WD}, \po{expr/WD}: well-definedness of an expression}
  \champ{Trigger.} Any expression or predicate $\mathit{expr}$ written in the model
  (axiom, theorem, invariant, guard, action, variant, witness).
  \champ{PO produced.}
  \[ \AX \vdash \WD(\mathit{expr}) \]
  (in a machine, the hypotheses also include the invariant and, where relevant,
  the guards; see R8).
  \champ{Computing $\WD$.} $\WD$ is defined by induction on the syntax, using three
  tables (predicates, unary/binary expressions, other expressions). The
  significant cases are:
  \begin{center}
  \begin{tabular}{>{$}l<{$} >{$}l<{$}}
    \toprule
    e & \WD(e)\\ \midrule
    P\wedge Q,\ P\limp Q & \WD(P)\wedge(P\limp \WD(Q))\\
    P\vee Q & \WD(P)\wedge(P\vee \WD(Q))\\
    P\Leftrightarrow Q & \WD(P)\wedge \WD(Q)\\
    \neg P & \WD(P)\\
    \forall L\cdot P,\ \exists L\cdot P & \forall L\cdot \WD(P)\\
    \lambda P\cdot Q\mid E & \forall \mathcal{F}_P\cdot \WD(Q)\wedge(Q\limp\WD(E))\\
    \bigcap L\cdot P\mid E & \forall L\cdot(\WD(P)\wedge(P\limp\WD(E)))\wedge\exists L\cdot P\\
    \bigcap E\mid P & \forall \mathcal{F}_E\cdot(\WD(P)\wedge(P\limp\WD(E)))\wedge\exists \mathcal{F}_E\cdot P\\
    F(E) & \WD(E)\wedge\WD(F)\wedge E\in\dom(F)\wedge F\in S\pfun T\\
    E / F & \WD(E)\wedge\WD(F)\wedge F\neq 0\\
    E \bmod F & \WD(E)\wedge\WD(F)\wedge 0\le E\wedge 0< F\\
    E^{F} & \WD(E)\wedge 0\le E\wedge\WD(F)\wedge 0\le F\\
    \mathrm{card}(E) & \WD(E)\wedge \mathrm{finite}(E)\\
    \mathrm{inter}(E) & \WD(E)\wedge E\neq\varnothing\\
    \min(E) & \WD(E)\wedge E\neq\varnothing\wedge \exists v\cdot\forall u\cdot u\in E\limp v\le u\\
    \max(E) & \WD(E)\wedge E\neq\varnothing\wedge \exists v\cdot\forall u\cdot u\in E\limp v\ge u\\
    E\ \mathsf{op}\ F \text{ (relations, $+$, $\cup$, \dots)} & \WD(E)\wedge\WD(F)\\
    \text{identifiers, literals, } \nat, \intg, \varnothing & \top\\
    \bottomrule
  \end{tabular}
  \end{center}
  In the $\lambda$ row, $P$ is the pattern (the bound variables) and $Q$ the
  predicate; $\mathcal{F}_X$ denotes the variables bound by $X$.
  \champ{Principle.} Connectives are evaluated \emph{from left to right and lazily}:
  in $P\wedge Q$, the well-definedness of $Q$ may assume $P$. This is why the order
  in which axioms, invariants and guards are written affects the POs produced
  (see R8).
  \champ{Example from the text.} In this example, $s$ is the sum function of the
  context, not the list of sets:
  $\WD(\forall k\cdot k\in\nat\limp 2*s(k)=k*k+k)
   = \forall k\cdot k\in\nat\limp k\in\dom(s)\wedge s\in\intg\pfun\intg$.
\end{regle}

\begin{regle}{R2 \quad \po{th/THM}: context theorem}
  \champ{Trigger.} A theorem labelled $th$ in a context (or a property on $s,c$
  placed among the theorems of a machine).
  \champ{PO produced.}
  \[ \AX,\ \Th(s,c) \vdash th \]
  where $\Th(s,c)$ denotes the theorems of the context that precede $th$.
  \champ{Justification.} The \sff{THEOREMS} clause contains consequences of the
  axiomatic theory: they must be \emph{proved}, not \emph{assumed}. The consistency
  of the theory itself is not checked by any PO; the chapter recommends checking it,
  for instance with ProB.
\end{regle}

% ==================================================================
\section{Rules for an abstract machine}

\begin{regle}{R3 \quad \po{th$_i$/THM}: safety property (machine theorem)}
  \champ{Trigger.} A theorem $th_i : \mathit{SAFE}_i(s,c,x)$ of the machine, $i\in 1..h$.
  \champ{Derivation.} $\AX\vdash I\limp \bigwedge_i \mathit{SAFE}_i$
  $\xrightarrow{\text{S1}}$ $\AX, I\vdash \bigwedge_i \mathit{SAFE}_i$
  $\xrightarrow{\text{S2}}$ one PO per $i$, to which the theorems already
  established are added.
  \champ{PO produced (for each $i$).}
  \[ \AX,\ \Th_i(s,c),\ I(s,c,x) \vdash \mathit{SAFE}_i(s,c,x) \]
  \champ{Role.} It separates the \emph{inductive} properties (invariants, which
  require an induction step) from the safety properties that \emph{follow} from
  them (theorems). In the \sff{SAFETY} example, $x\le 0$ is not inductive, but it
  follows from the inductive invariant $x=-1$.
\end{regle}

\begin{regle}{R4 \quad \po{INITIALISATION/inv$_i$/INV}: establishing the invariant}
  \champ{Trigger.} The \sff{INITIALISATION} event, written
  $x :\mid \mathit{initBAP}(s,c,x')$, and each invariant $inv_i$.
  \champ{Derivation.} $\AX\vdash \mathit{Init}\limp I$ $\xrightarrow{\text{S1}}$
  $\AX,\mathit{Init}\vdash I$ $\xrightarrow{\text{S2}}$ one PO per $inv_i$, stated by
  Rodin with the after-values $x'$.
 
  \[ \AX,\ \mathit{initBAP}(s,c,x') \vdash I_i(s,c,x') \]
  \champ{Remark.} The initialisation has no before-state: neither the invariant nor
  a typing hypothesis on $x$ appears as a hypothesis.
\end{regle}

\begin{regle}{R5 \quad \po{INITIALISATION/act$_i$/FIS}: feasibility of the initialisation}
  \champ{Trigger.} Each \emph{non-deterministic} action
  $act_i : x_i :\mid \mathit{initBAP}_i(s,c,x_i')$ of the initialisation, the $x_i$
  forming a \emph{partition} of $x$.
  \champ{Derivation.} $\AX\vdash\exists x'\cdot\mathit{initBAP}(s,c,x')$ with
  $\mathit{initBAP}\equiv\bigwedge_i\mathit{initBAP}_i$. Since the $x_i$ are
  disjoint, the existential distributes:
  $\exists x'\cdot\bigwedge_i \mathit{initBAP}_i(x_i') \equiv \bigwedge_i \exists x_i'\cdot \mathit{initBAP}_i(x_i')$,
  and S2 is then applied.
  \champ{PO produced.}
  \[ \AX \vdash \exists x_i'\cdot \mathit{initBAP}_i(s,c,x_i') \]
  \champ{Role.} It guarantees that the model has at least one initial state. When
  the action gives no explicit value (for instance
  $acc:\mid acc'\in P\to B\wedge acc'\subseteq a$), a witness constant $acc0$ is
  introduced in the context. The chapter points out that this only moves the
  difficulty: the axiom on $acc0$ is consistent only if $a$ contains a total
  function from $P$ to $B$. For a deterministic assignment $x:=e$, no
  \po{FIS} PO is produced; only the \po{WD} PO for $e$ remains.
\end{regle}

\begin{regle}{R6 \quad \po{e/inv$_i$/INV}: invariant preservation by an event}
  \champ{Trigger.} Each event \sff{e} (other than the initialisation) and each
  invariant $inv_i$.
  \champ{Derivation.} From (INV): $\AX\vdash I\wedge \BA(\sff{e})\limp I'$, unfold
  $\BA$, eliminate $\exists u$ in the hypothesis (S3), apply S1, then S2.
  \champ{PO produced.}
  \[ \AX,\ I(s,c,x),\ G(u,s,c,x),\ \mathit{BAP}(u,s,c,x,x') \vdash I_i(s,c,x') \]
  \champ{Remark.} The \emph{whole} invariant is a hypothesis; only the conclusion is
  split. In the \sff{SAFETY} example, the PO for $x=-1$ holds because the
  hypotheses are inconsistent ($x=-1$, $x\ge 0$). With $x\le 0$, one would obtain
  the false PO $\dots\vdash 1\le 0$, which shows that $x\le 0$ is not inductive.
\end{regle}

\begin{regle}{R7 \quad \po{e/act$_i$/FIS}: feasibility of an event}
  \champ{Trigger.} Each \emph{non-deterministic} action
  $act_i : x_i :\mid \mathit{BAP}_i(u,s,c,x,x_i')$ (or $x_i :\in E$) of the event \sff{e}.
  
  \[ \AX,\ I(s,c,x),\ G(u,s,c,x) \vdash \exists x_i'\cdot \mathit{BAP}_i(u,s,c,x,x_i') \]
  \champ{Role.} It corresponds to the $\Fis(\sff{e})$ part of the \emph{wp}
  condition: when the guard holds, the event can actually produce a next state. As
  for the initialisation, the existential distributes over the actions, which act
  on disjoint variables, hence one PO per action. For $x:=e$, the existence of
  $x'=e$ is trivial and no \po{FIS} PO is produced.
\end{regle}
\begin{rodin}
  Box \emph{feasibility in Rodin} of the chapter:
  \begin{itemize}
    \item \po{FIS} is generated only for non-deterministic actions, $x:\in E$ and
      $x:\mid P(x,x')$; for $x:\in E$ it reduces to $E\neq\varnothing$; an action
      $x:=e$ produces only \po{WD};
    \item feasibility is split action by action; this is sound because the actions
      modify disjoint lists of variables;
    \item limitation: the feasibility of an action is proved in isolation, and an
      unprovable \po{FIS} PO does not say whether to strengthen the guard or weaken
      the action.
  \end{itemize}
\end{rodin}

\begin{regle}{R8 \quad \po{e/grd$_i$/WD}, \po{e/act$_i$/WD} and \po{INITIALISATION/act$_i$/WD}: well-definedness of guards and actions}
  \champ{Trigger.} Each guard and each action of an event, and each action of the
  initialisation.
  
  \begin{align*}
    &\AX,\ I(s,c,x),\ G_1,\dots,G_{i-1} \vdash \WD(G_i(u,s,c,x)) && (\po{e/grd$_i$/WD})\\
    &\AX,\ I(s,c,x),\ G(u,s,c,x) \vdash \WD(\mathit{BAP}_i(u,s,c,x,x_i')) && (\po{e/act$_i$/WD})\\
    &\AX \vdash \WD(\mathit{initBAP}_i(s,c,x_i')) && (\po{INITIALISATION/act$_i$/WD})
  \end{align*}
  \champ{Role.} They rule out ``silly'' expressions (in Lamport's sense), such as a
  division by zero. The guards that precede $G_i$ are hypotheses: the list of
  guards is read as the conjunction $G_1\wedge\dots\wedge G_k$, to which the rule
  $\WD(P\wedge Q)=\WD(P)\wedge(P\limp\WD(Q))$ is applied. This allows one to write
  the two guards $grd_1 : x\in\dom(f)$ and $grd_2 : f(x)>0$.
\end{regle}
\begin{rodin}
  Box \emph{well-definedness in Rodin} of the chapter:
  \begin{itemize}
    \item the same left-to-right reading applies to every labelled list: an axiom,
      an invariant or a guard is proved well defined with those that precede it,
      and an action with all the guards;
    \item order matters: exchanging $grd_1$ and $grd_2$ above yields
      $\AX, I\vdash x\in\dom(f)$, which is in general not provable;
    \item within a single predicate, $\wedge$ is not commutative for WD:
      $x\in\dom(f)\wedge f(x)>0$ is well defined, $f(x)>0\wedge x\in\dom(f)$ is not;
    \item Rodin discharges most \po{WD} POs automatically: an undischarged \po{WD} PO
      usually points to a missing or misplaced hypothesis.
  \end{itemize}
\end{rodin}

\subsection{Event status and variants}\label{sec:statut}
Each event has a status: \emph{ordinary}, \emph{convergent} or
\emph{anticipated}. The text justifies the variant POs by the following liveness
rule. If $R_{n+1}\leadsto R_n$ for all $n$, if $P\wedge I\limp\exists n\cdot R_n$
and if $R_0\limp Q$, then $P\leadsto Q$, with
$R_n \widehat{=} I\wedge \mathit{varexp}\le n$ (the variant may decrease
by more than one unit). The POs ensure that a convergent event strictly decreases a
well-founded quantity, hence that convergent events cannot occur forever. To
conclude $P\leadsto Q$, a convergent event must also remain enabled as long as $Q$
does not hold: this follows from deadlock freedom or from a fairness assumption.

\begin{regle}{R9 \quad \po{e/NAT} and \po{e/FIN}: the variant is well founded}
  \champ{Trigger.} A \emph{convergent} (or anticipated) event \sff{e}. A machine has
  a single variant (clause \sff{variant}).
  \champ{PO produced (numeric variant).}
  \[ \AX,\ I(s,c,x),\ G_e(u,s,c,x) \vdash \mathit{varexp}(s,c,x)\in\nat \]
  \champ{PO produced (set-valued variant).}
  \[ \AX,\ I(s,c,x),\ G_e(u,s,c,x) \vdash \mathrm{finite}(\mathit{varexp}(s,c,x)) \]
\end{regle}

\begin{regle}{R10 \quad \po{e/VAR}: decrease of the variant}
  \champ{Trigger.} A \emph{convergent} event \sff{e}.
  \champ{PO produced (numeric variant).}
  \[ \AX,\ I,\ G_e,\ \mathit{BAP}(u,s,c,x,x') \vdash \mathit{varexp}(s,c,x') < \mathit{varexp}(s,c,x) \]
  \champ{PO produced (set-valued variant).}
  \[ \AX,\ I,\ G_e,\ \mathit{BAP}(u,s,c,x,x') \vdash \mathit{varexp}(s,c,x') \subset \mathit{varexp}(s,c,x) \]
  \champ{For an anticipated event.} The variant is only required not to increase:
  $\le$ (resp.\ $\subseteq$). See R18.
\end{regle}

% ==================================================================
\section{Rules for refinement}

\subsection{Reference definition and derivation}\label{sec:defraf}
\begin{regle}{Definition (refinement of two events, \emph{wp} style)}
  $\mathsf{c}$ refines $\mathsf{a}$ with respect to $x, I, y, J$ if
  \[ \AX \vdash I(s,c,x)\wedge J(s,c,x,y) \limp [\mathsf{c}]\,\neg[\mathsf{a}]\,\neg J(s,c,x,y). \]
\end{regle}
This is the definition of operation refinement in B, transposed to events. The
chapter transforms it in seven steps, all of them equivalences: unfolding
$[\mathsf{a}]$ (normalised form with $\exists u$), simplifying the negation,
unfolding $[\mathsf{c}]$ (with $\exists v$), eliminating $\forall y'$ (free
variable), implication introduction (S1), eliminating $\exists v$ in the hypothesis
(S3). With $\Gamma_0 \widehat{=} \AX, I(s,c,x), J(s,c,x,y)$, one obtains:
\begin{multline*}
  \text{(7)}\qquad \Gamma_0,\ H(v,s,c,y),\ \mathit{CBAP}(v,s,c,y,y') \\
  \vdash \exists x'\cdot\exists u\cdot\big(G(u,s,c,x)\wedge\mathit{ABAP}(u,s,c,x,x')\wedge J(s,c,x',y')\big).
\end{multline*}
The definition does not guarantee the feasibility of $\mathsf{c}$: it is proved
separately (\po{c/act$_i$/FIS}, R16).

\begin{regle}{Property (event refinement, version II): GRD and SIM}
  $\mathsf{c}$ refines $\mathsf{a}$ \emph{if, and only if,} (7) holds. Consequently,
  if $\mathsf{c}$ refines $\mathsf{a}$, then
  \begin{align*}
    \text{(GRD)}\quad & \AX, I, J, H(v,s,c,y) \vdash \exists u\cdot G(u,s,c,x)\\
    \text{(SIM)}\quad & \AX, I, J, H(v,s,c,y), \mathit{CBAP}(v,s,c,y,y')\\
      & \qquad \vdash \exists x'\cdot\exists u\cdot \mathit{ABAP}(u,s,c,x,x')\wedge J(s,c,x',y')
  \end{align*}
  The converse does not hold in general: GRD and SIM are \emph{necessary} but not
  \emph{sufficient}.
\end{regle}
Going from (7) to GRD and SIM separates two existentials that referred to the
\emph{same} value of $u$. Counterexample: $u\in\{0,1\}$, $G\equiv u=0$,
$\mathit{ABAP}\equiv u=1\wedge x'=x+1$, $J\equiv\mathit{true}$; GRD holds with
$u=0$, SIM with $u=1$, but (7) is false. In GRD, the hypothesis $\mathit{CBAP}$ is
dropped: the abstract guard does not depend on $y'$; this is the form used by
Rodin. The chapter's summary list (property \texttt{pro:poraf}) now
follows this form for \po{GRD}, \po{GRD-WIT} and \po{WFIS-P}.

The witnesses restore the link. Rule S4 is applied with the witnesses declared in
the \sff{witness} clause of the concrete event: $u : \mathit{WP}(u,s,c,v,y)$ and
$x' : \mathit{WV}(v,s,c,y',x')$.

\begin{regle}{Property (refinement with witnesses, version III)}
  If the POs \po{GRD-WIT} (R11), \po{SIM-WIT} (R12 and R13), \po{WFIS-P} and
  \po{WFIS-V} (R14) are valid, then $\mathsf{c}$ refines $\mathsf{a}$.
  \champ{Proof (chapter).} \po{WFIS-P} gives $u_0$ with $\mathit{WP}(u_0)$,
  \po{WFIS-V} gives $x'_0$ with $\mathit{WV}(x'_0)$; \po{GRD-WIT} gives $G(u_0)$
  and \po{SIM-WIT} gives $\mathit{ABAP}(u_0,x'_0)\wedge J(x'_0,y')$, hence (7). The
  key point is that \po{GRD-WIT} and \po{SIM-WIT} share the hypothesis
  $\mathit{WP}(u,s,c,v,y)$: the same value of $u$ is used for both. The witnesses
  therefore make the decomposition sound; they are not merely an aid for the
  prover.
\end{regle}
\textbf{When is a witness required?} For every abstract parameter that disappears
in the concrete event, and for every abstract variable that disappears and that the
abstract event modifies non-deterministically. Elsewhere, witnesses are optional
but help the prover (see \texttt{temoins-obligatoires.pdf}).

\textbf{Machine refinement.} \sff{CM} refines \sff{AM} if the concrete
initialisation refines the abstract one (R15) and if every event of \sff{CM}
refines an event of \sff{AM}, the new events refining \sff{skip}.

\subsection{Rules generated for a concrete event}

\begin{regle}{R11 \quad \po{c/grd/GRD}: guard strengthening}
  \champ{Trigger.} A concrete event $\mathsf{c}$ \sff{refines} $\mathsf{a}$, and each
  abstract guard $grd : G$.
  \champ{PO produced (with a witness).}
  \[ \AX,\ I(s,c,x),\ J(s,c,x,y),\ H(v,s,c,y),\ \mathit{WP}(u,s,c,v,y) \vdash G(u,s,c,x) \]
  \champ{Role.} When the concrete event is observable, so is the abstract event:
  the concrete model shows no new behaviour.
\end{regle}

\begin{regle}{R12 \quad \po{c/jnv$_i$/INV}: preservation of the gluing invariant}
  \champ{Trigger.} $\mathsf{c}$ and each invariant $jnv_i : J_i$ of the concrete
  machine.
  \champ{Derivation.} The conclusion of SIM-WIT, $\mathit{ABAP}\wedge J'$, is split
  by S2 into an INV part and a SIM part.
  
  \begin{multline*} \AX,\ I,\ J,\ H(v,s,c,y),\ \mathit{CBAP}(v,s,c,y,y'),\\ \mathit{WP}(u,s,c,v,y),\ \mathit{WV}(v,s,c,y',x') \vdash J_i(s,c,x',y') \end{multline*}
\end{regle}

\begin{regle}{R13 \quad \po{c/act/SIM}: simulation of the abstract action}
  \champ{Trigger.} $\mathsf{c}$ and each abstract action $act : x :\mid \mathit{ABAP}$
  of $\mathsf{a}$.
  \champ{PO produced.}
  \begin{multline*} \AX,\ I,\ J,\ H(v,s,c,y),\ \mathit{CBAP}(v,s,c,y,y'),\\ \mathit{WP}(u,s,c,v,y),\ \mathit{WV}(v,s,c,y',x') \vdash \mathit{ABAP}(u,s,c,x,x') \end{multline*}
  \champ{Role.} Every concrete transition corresponds to an abstract transition.
  The abstract invariant $I$ is not \emph{proved again}: it is inherited, since
  $\mathsf{a}$ already preserves it.
\end{regle}

\begin{regle}{R14 \quad \po{c/u/WFIS}, \po{c/x'/WFIS}, \po{c/u/WWD} and \po{c/x'/WWD}: feasibility and well-definedness of witnesses}\label{sec:wfis}
  \champ{Trigger.} Each witness declared in the \sff{witness} clause.
  \champ{Feasibility POs (non-deterministic witnesses).}
  \begin{align*}
    &\AX,\ I,\ J,\ H(v,s,c,y) \vdash \exists u\cdot \mathit{WP}(u,s,c,v,y)\\
    &\AX,\ I,\ J,\ H(v,s,c,y),\ \mathit{CBAP}(v,s,c,y,y') \vdash \exists x'\cdot \mathit{WV}(v,s,c,y',x')
  \end{align*}
 
  \begin{align*}
    &\AX,\ I,\ J,\ H(v,s,c,y) \vdash \WD(\mathit{WP}(u,s,c,v,y))\\
    &\AX,\ I,\ J,\ H(v,s,c,y),\ \mathit{CBAP}(v,s,c,y,y') \vdash \WD(\mathit{WV}(v,s,c,y',x'))
  \end{align*}
  \champ{Role.} The \po{WFIS} POs make the use of witnesses sound: together with
  R11 to R13, they imply (7) (rule S4). The \po{WWD} POs treat a witness like any
  other formal text written by the modeller.
\end{regle}
\begin{rodin}
  Box \emph{witnesses in Rodin} of the chapter.
  \begin{itemize}
    \item Rodin checks statically that a witness is given for every abstract
      parameter that disappears and for every disappearing abstract variable that
      the abstract event modifies non-deterministically; a missing witness is
      reported as an error, before any proof;
    \item \po{WFIS} is generated only for non-deterministic witnesses: for a
      deterministic witness $u=E$ or $x'=E$, Rodin substitutes $E$ for the variable
      in the other POs (one-point rule); only \po{WWD} remains;
    \item \po{WWD} is generated for all witnesses;
    \item limitation: the tool does not help in finding a witness; a wrong choice
      shows up only as an unprovable \po{GRD}, \po{SIM}, \po{INV} or \po{WFIS} PO,
      and the modeller must then revise the witness or the gluing invariant.
  \end{itemize}
\end{rodin}

\begin{regle}{R15 \quad Refined initialisation: \po{INITIALISATION/jnv$_i$/INV}, \po{act/SIM}, \po{act$_i$/FIS}, \po{act$_i$/WD}, \po{x'/WFIS}, \po{x'/WWD}}
  \champ{Global condition.}
  $\AX, \mathit{CInit}(s,c,y') \vdash \exists x'\cdot \mathit{AInit}(s,c,x')\wedge J(s,c,x',y')$.
  \champ{With the witness $x' : \mathit{WV}(s,c,y',x')$ and S2.} The initialisation
  has no parameter: the witness does not depend on $v$.
  \begin{align*}
    &\AX,\ \mathit{CInit}(s,c,y'),\ \mathit{WV}(s,c,y',x') \vdash J_i(s,c,x',y') && (\po{INITIALISATION/jnv$_i$/INV})\\
    &\AX,\ \mathit{CInit}(s,c,y'),\ \mathit{WV}(s,c,y',x') \vdash \mathit{AInit}(s,c,x') && (\po{INITIALISATION/act/SIM})\\
    &\AX \vdash \exists y'\cdot \mathit{CInit}(s,c,y') && (\po{INITIALISATION/act/FIS})\\
    &\AX \vdash \WD(\mathit{CInit}_i(s,c,y_i')) && (\po{INITIALISATION/act$_i$/WD})\\
    &\AX,\ \mathit{CInit}(s,c,y') \vdash \exists x'\cdot \mathit{WV}(s,c,y',x') && (\po{INITIALISATION/x'/WFIS})\\
    &\AX,\ \mathit{CInit}(s,c,y') \vdash \WD(\mathit{WV}(s,c,y',x')) && (\po{INITIALISATION/x'/WWD})
  \end{align*}
  \po{x'/WFIS} is produced only when the witness is not deterministic.
\end{regle}

\begin{regle}{R16 \quad \po{c/act$_i$/FIS}, \po{c/grd$_i$/WD}, \po{c/act$_i$/WD}: feasibility and well-definedness of the concrete event}
  \champ{Trigger.} The guards and actions of the concrete event $\mathsf{c}$.
  
  with the abstract and gluing invariants as hypotheses:
  \begin{align*}
    &\AX,\ I,\ J,\ H(v,s,c,y) \vdash \exists y_i'\cdot \mathit{CBAP}_i(v,s,c,y,y_i') && (\po{c/act$_i$/FIS})\\
    &\AX,\ I,\ J,\ H_1,\dots,H_{i-1} \vdash \WD(H_i(v,s,c,y)) && (\po{c/grd$_i$/WD})\\
    &\AX,\ I,\ J,\ H(v,s,c,y) \vdash \WD(\mathit{CBAP}_i(v,s,c,y,y_i')) && (\po{c/act$_i$/WD})
  \end{align*}
  As in R7, \po{c/act$_i$/FIS} is produced only for non-deterministic actions.
\end{regle}

\begin{regle}{R17 \quad \po{th/THM}: theorem of the concrete machine}
  \[ \AX,\ \Th(s,c),\ I(s,c,x),\ J(s,c,x,y) \vdash \mathit{SAFE}(s,c,x,y) \]
\end{regle}

\begin{regle}{R18 \quad \po{c/NAT}, \po{c/FIN}, \po{c/VAR}: variants in a refinement}
  \champ{Trigger.} A concrete event with status convergent or anticipated; the
  variant refers only to the concrete variables $y$.
  \begin{align*}
    &\AX, I, J, H(v,s,c,y) \vdash \mathit{varexp}(s,c,y)\in\nat && (\po{c/NAT})\\
    &\AX, I, J, H(v,s,c,y) \vdash \mathrm{finite}(\mathit{varexp}(s,c,y)) && (\po{c/FIN}, \text{ set-valued variant})\\
    &\AX, I, J, H, \mathit{CBAP} \vdash \mathit{varexp}(s,c,y') < \mathit{varexp}(s,c,y) && (\text{convergent})\\
    &\AX, I, J, H, \mathit{CBAP} \vdash \mathit{varexp}(s,c,y') \le \mathit{varexp}(s,c,y) && (\text{anticipated})
  \end{align*}
  Rodin produces one or the other \po{c/VAR} PO depending on the status of the event;
  for a set-valued variant, $<$ and $\le$ become $\subset$ and $\subseteq$.
  \champ{Methodology.} A \emph{new} event implicitly refines \sff{skip}: its SIM PO
  then reduces to $x'=x$. It must be convergent so as not to take control forever.
  An anticipated event postpones this obligation to a later refinement, while
  guaranteeing that it does not increase the variant.
\end{regle}

\begin{regle}{R19 \quad \po{c/MRG}: event merging}
  \champ{Trigger.} A concrete event $\mathsf{c}$ that refines \emph{two} abstract
  events $\mathsf{a}_1$ and $\mathsf{a}_2$ (\sff{c(a1+a2)}).
  \champ{Application condition (syntactic).} The actions of $\mathsf{a}_1$,
  $\mathsf{a}_2$ and $\mathsf{c}$ must be \emph{syntactically identical}, and the merged
  events must have the same parameters.
  \champ{PO produced.}
  \[ \AX,\ I(s,c,x),\ J(s,c,x,y),\ H(v,s,c,y) \vdash G_1(u,s,c,x)\vee G_2(u,s,c,x) \]
  The gluing invariant $J$ links $y$ (in $H$) to $x$ (in $G_1,G_2$).
  \champ{Role.} The GRD PO is replaced by the strengthening of the
  \emph{disjunction} of the guards. The rule makes it possible to reduce the number
  of events, for instance to obtain a conditional. In the example of the maximum of
  two numbers (machines \sff{MRG0} to \sff{MRG3}), the PO becomes
  $done=\mathit{FALSE}\vdash(i<j\wedge done=\mathit{FALSE})\vee(i\ge j\wedge done=\mathit{FALSE})$,
  which reduces to $i<j\vee i\ge j$, hence to true.
\end{regle}

% ==================================================================
\section{Deadlock freedom (semantics)}

These POs are \emph{not} generated systematically by Rodin or by Atelier~B: the
modeller adds them as theorems, or checks them by model checking (ProB).

\begin{regle}{R20 \quad \po{m/DLF}: deadlock-free machine}
  \[ \AX,\ I(s,c,x) \vdash \bigvee_{\sff{e}\in E(\mathsf{AM})} \exists u\cdot G_{\sff{e}}(u,s,c,x) \]
\end{regle}
\begin{regle}{R21 \quad \po{r/DLF}: the refinement does not deadlock more often}
  \[ \AX,\ I,\ J,\ \bigvee_{\sff{e}\in E(\mathsf{AM})} \exists u\cdot G_{\sff{e}}(u,s,c,x)
     \vdash \bigvee_{\sff{e}\in E(\mathsf{CM})} \exists v\cdot H_{\sff{e}}(v,s,c,y) \]
\end{regle}

% ==================================================================
\section{Summary table}

\begin{longtable}{p{0.27\linewidth} p{0.27\linewidth} p{0.38\linewidth}}
\toprule
\textbf{Name} & \textbf{Trigger} & \textbf{Hypotheses $\vdash$ goal}\\ \midrule\endhead
\po{x/WD} & any expression & $\AX\,[,I,G] \vdash \WD(x)$\\
\po{th/THM} (context) & context theorem & $\AX,\Th\vdash th$\\
\po{th$_i$/THM} (machine) & machine theorem & $\AX,\Th_i,I\vdash \mathit{SAFE}_i$\\
\po{INITIALISATION/inv$_i$/INV} & invariant $\times$ init. & $\AX,\mathit{initBAP}'\vdash I_i'$\\
\po{INITIALISATION/act$_i$/FIS} & non-det. init. action & $\AX\vdash\exists x_i'\cdot\mathit{initBAP}_i$\\
\po{INITIALISATION/act$_i$/WD} & init. action & $\AX\vdash\WD(\mathit{initBAP}_i)$\\
\po{e/inv$_i$/INV} & invariant $\times$ event & $\AX,I,G,\mathit{BAP}\vdash I_i'$\\
\po{e/act$_i$/FIS} & non-det. action & $\AX,I,G\vdash\exists x_i'\cdot\mathit{BAP}_i$\\
\po{e/grd$_i$/WD} & guard & $\AX,I,G_1..G_{i-1}\vdash \WD(G_i)$\\
\po{e/act$_i$/WD} & action & $\AX,I,G\vdash \WD(\mathit{BAP}_i)$\\
\po{e/NAT} & convergent/anticipated event & $\AX,I,G\vdash \mathit{varexp}\in\nat$\\
\po{e/FIN} & same, set-valued variant & $\AX,I,G\vdash \mathrm{finite}(\mathit{varexp})$\\
\po{e/VAR} & convergent event & $\AX,I,G,\mathit{BAP}\vdash \mathit{varexp}'<\mathit{varexp}$ (or $\subset$)\\
\po{c/grd/GRD} & abstract guard $\times$ refinement & $\AX,I,J,H,\mathit{WP}\vdash G$\\
\po{c/act/SIM} & abstract action $\times$ refinement & $\AX,I,J,H,\mathit{CBAP},\mathit{WP},\mathit{WV}\vdash \mathit{ABAP}$\\
\po{c/jnv$_i$/INV} & gluing invariant & $\AX,I,J,H,\mathit{CBAP},\mathit{WP},\mathit{WV}\vdash J_i'$\\
\po{c/u/WFIS} & non-det. parameter witness & $\AX,I,J,H\vdash\exists u\cdot \mathit{WP}$\\
\po{c/x'/WFIS} & non-det. variable witness & $\AX,I,J,H,\mathit{CBAP}\vdash\exists x'\cdot \mathit{WV}$\\
\po{c/u/WWD}, \po{c/x'/WWD} & any witness & same hypotheses $\vdash\WD(\mathit{WP})$, $\WD(\mathit{WV})$\\
\po{c/act$_i$/FIS} & non-det. concrete action & $\AX,I,J,H\vdash\exists y_i'\cdot\mathit{CBAP}_i$\\
\po{c/grd$_i$/WD}, \po{c/act$_i$/WD} & concrete guard, action & $\AX,I,J,H_1..H_{i-1}\vdash\WD(H_i)$; $\AX,I,J,H\vdash\WD(\mathit{CBAP}_i)$\\
\po{INITIALISATION/jnv$_i$/INV} & refined init. & $\AX,\mathit{CInit}',\mathit{WV}\vdash J_i(x',y')$\\
\po{INITIALISATION/act/SIM} & refined init. & $\AX,\mathit{CInit}',\mathit{WV}\vdash \mathit{AInit}(x')$\\
\po{INITIALISATION/x'/WFIS}, \po{x'/WWD} & init. witness & $\AX,\mathit{CInit}'\vdash\exists x'\cdot\mathit{WV}$; $\vdash\WD(\mathit{WV})$\\
\po{th/THM} (refinement) & concrete theorem & $\AX,\Th,I,J\vdash\mathit{SAFE}$\\
\po{c/NAT}, \po{c/FIN} & convergent / anticipated & $\AX,I,J,H\vdash \mathit{varexp}(y)\in\nat$ / finite\\
\po{c/VAR} & convergent / anticipated & $\dots\vdash \mathit{varexp}(y')<\mathit{varexp}(y)$ / $\le$\\
\po{c/MRG} & merging $\mathsf{a}_1+\mathsf{a}_2$ & $\AX,I,J,H\vdash G_1\vee G_2$\\
\po{m/DLF}, \po{r/DLF} & optional & disjunction of the guards\\
\bottomrule
\end{longtable}


\bibliography{refs}
\end{document}